% 02_layout.tex
% Geometría, Encabezados, Títulos y Metadatos PDF (ISO 15289)

% 1. Configuración de Márgenes Exactos
\geometry{
    a4paper,
    inner=2.5cm,       % Margen izquierdo para archivo físico
    outer=2.5cm,       % Margen derecho
    top=2.5cm,       % Distancia superior
    bottom=2.5cm,    % Distancia inferior
    includehead,     % El encabezado va dentro del margen
    includefoot,     % El pie de página va dentro del margen
    headheight=28pt, 
    headsep=0.4cm,   
    footskip=1.0cm   
}
\onehalfspacing      % Interlineado 1.5 para legibilidad en lectura larga

% 2. Jerarquía de Títulos (Especificaciones de Diseño)
% Portada de capítulo estilo "libro": el título del capítulo ocupa su propia
% plana y el contenido comienza en la página siguiente. Se logra con un
% \titleformat de tipo "display" que coloca un rótulo "CAPÍTULO N" y el título
% en gran tamaño. NO se dibuja regla horizontal; en su lugar, cada capítulo
% declara un breve resumen con \resumenCapitulo{...}, que cierra la plana.
\titleformat{\chapter}[display]
  {\sffamily\bfseries\color{colorPrimario}}
  {\filright\fontsize{13}{15}\selectfont\color{grisISO}\MakeUppercase{Capítulo \thechapter}}
  {0.35cm}
  {\fontsize{25}{29}\selectfont\filright}
  []
\titlespacing*{\chapter}{0pt}{55pt}{24pt}

% \resumenCapitulo{texto}: sinopsis breve bajo el título del capítulo, en la
% misma plana. Cierra con \clearpage para que el contenido empiece en la página
% siguiente (estilo libro). Sustituye a la antigua regla horizontal de color.
\newcommand{\resumenCapitulo}[1]{%
  \vspace{0.5cm}%
  {\color{colorAcento}\rule{2.5cm}{2pt}}\par
  \vspace{0.45cm}%
  {\normalfont\sffamily\fontsize{11}{15}\selectfont\color{grisISO!95!black}%
   \itshape #1\par}%
  \clearpage}

% Subtítulos: Helvetica, 14 pt, Negrita
\titleformat{\section}[hang]
  {\sffamily\bfseries\fontsize{14}{16.8}\selectfont\color{colorPrimario}}
  {\thesection.}
  {1em}{}

% 3. Identidad Corporativa: Encabezado y Pie (Trazabilidad)

% Variables de trazabilidad documental (ISO 15289). Se redefinen desde main.tex.
\newcommand{\currentDocID}{MU-PERF-000}
\newcommand{\currentDocEstado}{Borrador}
\newcommand{\setDocID}[1]{\renewcommand{\currentDocID}{#1}}
\newcommand{\setDocEstado}[1]{\renewcommand{\currentDocEstado}{#1}}

\pagestyle{fancy}
\fancyhf{}

% Redefinimos las marcas para capturar solo el nombre limpio del capítulo/sección
\renewcommand{\chaptermark}[1]{\markboth{#1}{}}
\renewcommand{\sectionmark}[1]{\markright{#1}}

% Encabezado: Logo (Izq) | Capítulo + Estado (Der)
\fancyhead[L]{\includegraphics[height=0.5cm, keepaspectratio]{assets/img/bytebusters/logo_horizontal.png}}

% \leftmark inyecta automáticamente el nombre del capítulo actual
\fancyhead[R]{\sffamily\footnotesize\textcolor{colorPrimario}{\textbf{\leftmark\ | Estado: \currentDocEstado}}}

% Pie de Página: ID Documental y Confidencialidad (Izq) | Paginación (Der)
\fancyfoot[L]{\sffamily\footnotesize\textcolor{grisISO}{\currentDocID\ | Confidencial - Byte Busters S.R.L.}}
\fancyfoot[R]{\sffamily\footnotesize\textcolor{grisISO}{Página \thepage\ de \pageref{LastPage}}}

\renewcommand{\headrulewidth}{0.4pt}
\renewcommand{\footrulewidth}{0.4pt}
\renewcommand{\headrule}{\hbox to\headwidth{\color{grisISO}\leaders\hrule height \headrulewidth\hfill}}
\renewcommand{\footrule}{\hbox to\headwidth{\color{grisISO}\leaders\hrule height \footrulewidth\hfill}}

% 4. Configuración de Metadatos y Enlaces
\hypersetup{
    colorlinks=true,
    linkcolor=colorPrimario,
    urlcolor=colorPrimario,
    pdftitle={Manual del Sistema para Perfiles - EthosHub (ISO/IEC/IEEE 26514:2022)},
    pdfauthor={Byte Busters S.R.L.}
}